% Mean-Variance Portfolio Optimization - Numerical Optimization course report
% Compile from the Docs directory with:
%   pdflatex -jobname=PortfolioOptimization_Report PortfolioOptimization_LatexSourceCode.tex
%   pdflatex -jobname=PortfolioOptimization_Report PortfolioOptimization_LatexSourceCode.tex

\documentclass[conference]{IEEEtran}

\usepackage{amsmath,amssymb}
\usepackage{booktabs}
\usepackage{cite}
\usepackage{graphicx}
\usepackage{url}
\usepackage{balance}

\newcommand{\vect}[1]{\boldsymbol{#1}}
\newcommand{\mat}[1]{\boldsymbol{#1}}
\newcommand{\norm}[1]{\left\lVert #1 \right\rVert}
\newcommand{\imageorplaceholder}[3]{%
    \IfFileExists{#1}{%
        \includegraphics[width=#2,height=#3,keepaspectratio]{#1}%
    }{%
        \fbox{%
            \parbox[c][#3][c]{\dimexpr#2-2\fboxsep-2\fboxrule\relax}{%
                \centering
                \textbf{Screenshot placeholder}\\[4pt]
                \texttt{\detokenize{#1}}\\[4pt]
                Replace the placeholder by adding the final application screenshot
                to the \texttt{Images} directory.
            }%
        }%
    }%
}

\title{Mean-Variance Portfolio Optimization via Active-Set Quadratic Programming}

\author{%
\IEEEauthorblockN{Tarik Agi\'{c}}
\IEEEauthorblockA{Student ID: 20108\\
Numerical Optimization Course Project}
}

\begin{document}

\maketitle

\begin{abstract}
This paper presents a complete implementation of long-only mean-variance
portfolio optimization as a convex quadratic program. Periodic asset returns
are loaded from a validated comma-separated-value matrix and used to estimate
the sample mean-return vector and unbiased sample covariance matrix. For a
specified target return, the optimizer minimizes portfolio variance subject to
full investment, the target-return equality, and nonnegative asset weights.
The numerical core is a primal active-set method. It starts from an explicitly
constructed feasible portfolio, solves an equality-constrained
Karush--Kuhn--Tucker system at every iteration, truncates a step when a bound
becomes active, and removes an active bound when its multiplier violates the
optimality conditions. The implementation uses the natID framework for the
graphical interface and for two interchangeable linear-system backends:
\texttt{dense::DblMatrix} and \texttt{sparse::DblSolver} with LU
factorization. Repeated solutions from the global minimum-variance return to
the largest feasible asset return form the efficient frontier. The application
also visualizes individual assets, the selected portfolio, changes of the
active set, and portfolio weights across target returns. A second bundled sample 
applies the same pipeline to 72 monthly observations for five market-traded ETFs 
derived from public adjusted-close data. On the bundled five-asset, 24-period data 
set, the global minimum-variance portfolio achieves a periodic return of $0.726733\%$ 
at $0.134594\%$ risk. At a $1\%$ target, the solver converges in six iterations 
with two active bounds and $0.281134\%$ risk. Dense and sparse solutions agree
to approximately machine precision. Automated tests additionally cover analytical 
cases, feasibility and KKT conditions, invalid inputs, frontier construction, and 
two 20-asset systems.
\end{abstract}

\begin{IEEEkeywords}
active-set method, efficient frontier, mean-variance optimization, quadratic
programming, portfolio optimization, sparse linear systems
\end{IEEEkeywords}

\section{Introduction}

Mean-variance portfolio selection expresses the allocation of capital as a
trade-off between expected return and risk. Markowitz's formulation established
that portfolio risk depends not only on the variances of individual assets but
also on their covariances \cite{markowitz1952}. Consequently, combining assets
whose returns do not move identically can yield a portfolio with lower
variance than any allocation derived from individual volatilities alone.

This project treats portfolio construction as a numerical optimization
problem rather than as a financial forecasting exercise. The inputs are a
matrix of periodic asset returns, and the outputs are optimal portfolio
weights for one or more prescribed returns. Short selling is excluded, every
unit of available capital must be invested, and no transaction costs, taxes,
or risk-free asset are modeled. Under these assumptions the problem has a
quadratic objective, two linear equalities, and bound constraints. Because a
sample covariance matrix is positive semidefinite, the resulting program is
convex; any point satisfying the Karush--Kuhn--Tucker (KKT) conditions is a
global solution.

The main numerical challenge is the changing set of zero-weight assets. An
asset with $w_i=0$ has an active lower bound. As the requested return changes,
assets enter and leave this active set, producing different allocation regimes
along the efficient frontier. Active-set quadratic-programming methods are
well suited to this structure because they explicitly identify the binding
constraints and solve a sequence of equality-constrained subproblems
\cite{nocedal2006,goldfarb1983}.

The completed application contains four connected stages. First, it validates
and loads return observations. Second, it estimates the mean vector and
covariance matrix. Third, it solves either a global minimum-variance problem or
a target-return problem by an active-set method. Finally, it constructs and
visualizes the efficient frontier and the corresponding weight paths. The
software is written in C++20 and uses natID both for its native graphical
interface and for dense and sparse matrix solvers \cite{natid}. A common
matrix-backend interface permits the same active-set algorithm to be executed
with either solver, making numerical agreement directly measurable.

The contribution of the project is therefore not a new financial model, but a
transparent end-to-end realization of the model: data validation, statistical
estimation, constrained optimization, backend comparison, visualization,
export, and automated numerical verification are implemented in one program.

\section{Problem Formulation}

\subsection{Return Statistics}

Let $T$ denote the number of observation periods and $n$ the number of assets.
The input matrix $\mat{R}\in\mathbb{R}^{T\times n}$ stores the periodic return
$R_{ti}$ of asset $i$ in period $t$. For each asset, the estimator of expected
periodic return is
\begin{equation}
    \mu_i = \frac{1}{T}\sum_{t=1}^{T} R_{ti}.
    \label{eq:mean}
\end{equation}
The implementation uses the unbiased sample covariance estimator
\begin{equation}
    \Sigma_{ij} = \frac{1}{T-1}\sum_{t=1}^{T}
    (R_{ti}-\mu_i)(R_{tj}-\mu_j).
    \label{eq:covariance}
\end{equation}
The standard deviation $\sigma_i=\sqrt{\Sigma_{ii}}$ is used when an individual
asset is drawn on the risk-return diagram. All quantities remain on the scale
of one observation period; the application does not silently annualize input
data.

For a portfolio weight vector $\vect{w}\in\mathbb{R}^{n}$, expected return and
variance are
\begin{align}
    \mu_p(\vect{w}) &= \vect{\mu}^{T}\vect{w}, \\
    \sigma_p^2(\vect{w}) &= \vect{w}^{T}\mat{\Sigma}\vect{w},
    \label{eq:metrics}
\end{align}
and the displayed portfolio risk is
$\sigma_p(\vect{w})=\sqrt{\sigma_p^2(\vect{w})}$.

\subsection{Long-Only Target-Return Program}

For a requested periodic return $r$, the application solves
\begin{align}
 \underset{\vect{w}}{\operatorname{minimize}}\quad
     & \frac{1}{2}\vect{w}^{T}\mat{\Sigma}\vect{w}
       \label{eq:qp}\\
 \operatorname{subject\ to}\quad
     & \vect{\mu}^{T}\vect{w}=r, \nonumber\\
     & \vect{1}^{T}\vect{w}=1, \nonumber\\
     & \vect{w}\geq \vect{0}. \nonumber
\end{align}
The factor $1/2$ does not change the minimizer, but makes the objective
gradient equal to $\mat{\Sigma}\vect{w}$. The budget equality enforces full
investment, while nonnegativity enforces a long-only portfolio. With only
these constraints, the feasible target interval is
\begin{equation}
    r\in\left[\min_i\mu_i,\ \max_i\mu_i\right].
    \label{eq:feasible-range}
\end{equation}
Targets outside this interval are rejected before a linear system is formed.

The global minimum-variance (GMV) problem removes the target-return equality
from \eqref{eq:qp}. Its solution is the left endpoint of the efficient branch
of the risk-return curve. The application first computes this portfolio and
then samples target values uniformly between its expected return and
$\max_i\mu_i$. Solving \eqref{eq:qp} at those targets yields the displayed
efficient frontier. Sampling begins at the GMV return rather than at
$\min_i\mu_i$, because lower-return minimum-variance portfolios belong to the
inefficient branch.

\section{Active-Set Numerical Method}

\subsection{Feasible Initialization}

A primal active-set method must begin with a feasible point. For an interior
target $r$, the implementation finds one asset with the largest mean below
$r$ and one with the smallest mean above $r$. If their returns are
$\mu_{\ell}<r<\mu_u$, a two-asset feasible portfolio is obtained from
\begin{align}
    w_{\ell} &= \frac{\mu_u-r}{\mu_u-\mu_{\ell}}, &
    w_u &= 1-w_{\ell},
    \label{eq:initial}
\end{align}
with all remaining weights equal to zero. Equation \eqref{eq:initial}
satisfies both equalities and nonnegativity by construction. At a boundary of
the feasible interval, the solver restricts the problem to all assets whose
means match that boundary and computes their minimum-variance allocation. The
GMV solve starts from equal weights $w_i=1/n$.

\subsection{Working Set and KKT Direction}

Let $\mathcal{W}_k$ be the set of zero-weight bounds treated as active at
iteration $k$, and let $\mat{E}_{\mathcal{W}_k}$ select their components. The
equality matrix and right-hand side are
\begin{equation}
 \mat{A}=\begin{bmatrix}\vect{1}^{T}\\\vect{\mu}^{T}\end{bmatrix},
 \qquad
 \vect{b}=\begin{bmatrix}1\\r\end{bmatrix}
 \label{eq:ab}
\end{equation}
for a target-return solve. The GMV version contains only the first row.

At a feasible iterate $\vect{w}_k$, the gradient is
$\vect{g}_k=\mat{H}\vect{w}_k$, where $\mat{H}$ is a regularized covariance
matrix described below. The search direction is obtained by solving the
equality-constrained quadratic subproblem. Its KKT equations are
\begin{equation}
\begin{bmatrix}
\mat{H} & \mat{A}^{T} & \mat{E}_{\mathcal{W}_k}^{T}\\
\mat{A} & \mat{0} & \mat{0}\\
\mat{E}_{\mathcal{W}_k} & \mat{0} & \mat{0}
\end{bmatrix}
\begin{bmatrix}
\vect{p}_k\\\vect{\lambda}_k\\\vect{\nu}_k
\end{bmatrix}
=
\begin{bmatrix}
-\vect{g}_k\\\vect{0}\\\vect{0}
\end{bmatrix}.
\label{eq:kkt}
\end{equation}
The last block forces $p_i=0$ for active assets, while
$\mat{A}\vect{p}_k=\vect{0}$ preserves the equalities. Direct factorization of
\eqref{eq:kkt} is appropriate for the moderate systems considered here;
nevertheless, the system is indefinite and can become ill-conditioned when
asset returns or covariance columns are nearly dependent \cite{golub2013}.

\subsection{Blocking Bounds and Multiplier Test}

If $\vect{p}_k$ is nonzero, the full step is accepted unless it would make a
free weight negative. The step length is
\begin{equation}
 \alpha_k=\min\left(1,
   \min_{\substack{i\notin\mathcal{W}_k\\p_{k,i}<0}}
   \frac{-w_{k,i}}{p_{k,i}}\right).
 \label{eq:step}
\end{equation}
The update $\vect{w}_{k+1}=\vect{w}_k+\alpha_k\vect{p}_k$ remains feasible.
When $\alpha_k<1$, the blocking weight is set exactly to zero and its bound is
added to the working set.

When $\norm{\vect{p}_k}_{\infty}$ is below the direction tolerance, the current
point minimizes the objective over the present working set. The multipliers of
the active bounds are then inspected. An active bound whose multiplier has the
wrong sign is removed, permitting that weight to increase in a later
iteration. If no multiplier violates the sign condition, the KKT conditions
for the full long-only problem are satisfied and the iteration terminates.
This add-or-remove logic is what allows the algorithm to follow different
active-set regions of the efficient frontier.

\subsection{Regularization and Stopping Checks}

The matrix used inside \eqref{eq:kkt} is symmetrized and slightly regularized:
\begin{align}
 \mat{H} &= \tfrac{1}{2}(\mat{\Sigma}+\mat{\Sigma}^{T})+\delta\mat{I},\\
 \delta &= \max\left(10^{-12},\rho\max(10^{-12},s)\right),\\
 s &= \frac{1}{n}\sum_{i=1}^{n}|\Sigma_{ii}|,
 \label{eq:regularization}
\end{align}
where the default relative regularization is $\rho=10^{-9}$. Regularization is
used only to stabilize the numerical systems. Reported variance and risk are
recomputed with the original sample covariance matrix.

The default maximum is 500 active-set iterations. The primal, direction, and
multiplier tolerances are respectively $10^{-9}$, $10^{-10}$, and $10^{-9}$.
After termination, the implementation independently checks equality residuals
and nonnegativity. A solve is reported as converged only when these checks
pass. This prevents a successful factorization from being confused with a
valid optimization result.

\section{Software Architecture and natID Integration}

The implementation separates data processing, numerical optimization, and
presentation. \texttt{CsvReader} parses the return matrix and reports the line
and column of invalid entries. It supports quoted CSV fields and an optional
observation-label column named \texttt{Date}, \texttt{Period}, \texttt{Time},
\texttt{Timestamp}, or \texttt{Observation}. Asset names must be nonempty and
unique, all numeric values must be finite, and the matrix must contain at least
two assets and two observations. \texttt{StatisticsCalculator} implements
\eqref{eq:mean} and \eqref{eq:covariance}. The structures
\texttt{ReturnData}, \texttt{Statistics}, \texttt{PortfolioSolution}, and
\texttt{EfficientFrontier} carry the resulting data between modules.

\texttt{ActiveSetQPSolver} implements the method of Section III and does not
depend on GUI controls. Linear algebra is delegated to
\texttt{MatrixBackend}, which exposes one interface for two professor-provided
natID implementations:
\begin{itemize}
    \item \texttt{dense::DblMatrix} stores the complete KKT matrix and solves
    it through the dense matrix API;
    \item \texttt{sparse::DblSolver} receives only nonzero triplets and uses
    general LU factorization with multi-pass Markowitz pivoting and the
    solver's own ordering strategy.
\end{itemize}
Although the mathematical KKT matrix is symmetric, the sparse path deliberately
uses the general LU mode. This avoids assuming positive definiteness, which an
indefinite saddle-point matrix does not possess. Sparse direct methods can
exploit structural zeros in larger systems, but their setup and ordering costs
mean that sparsity does not automatically imply lower runtime for a small
portfolio \cite{davis2006}.

\texttt{AppState} owns the loaded data and current results. It resets stale
solutions whenever a new data set is loaded, invokes the selected backend,
constructs summaries, compares both backends, and exports the frontier. The
export contains target and achieved returns, risk, variance, iteration count,
active-constraint names, and every asset weight at full double precision.

The natID interface contains separate views for return data, the central
Portfolio Explorer, the efficient frontier, selected allocation, and portfolio
weight paths. The Portfolio Explorer synchronizes a precise target-return input
with an interactive slider and presents the expected return, risk, variance,
and asset allocation of the selected solution. After the efficient frontier
has been constructed, three navigation landmarks---Minimum Risk, Balanced, and
Highest Return---select the first, middle, and final sampled frontier solutions,
respectively. These landmarks do not introduce a separate financial model or
recommendation criterion; they provide convenient access to solutions already
computed by the numerical optimizer.

Numerical information remains explicitly available in the same workspace,
including the selected natID matrix backend, active-set iteration count, active
zero-weight constraints, and dense--sparse comparison. Data-dependent controls
are disabled until a valid return matrix has been loaded, and the target-return
limits are adjusted to the feasible interval in
\eqref{eq:feasible-range}. The risk-return chart draws individual assets, the
efficient frontier, and a distinct marker for the selected portfolio. A
separate allocation view displays its weights, while the weight-path chart
plots $w_i(r)$ for every asset and marks sampled positions where the active set
changes. English and Bosnian resources use identical translation keys. A CSV
format panel documents the expected input directly inside the application.

\section{Experimental Setup and Verification}

\subsection{Bundled Demonstration Data}

The bundled deterministic data set contains 24 monthly observations for five
synthetic asset classes. Synthetic observations are appropriate for a
reproducible demonstration: every evaluator obtains identical statistics,
frontier points, and active-set transitions. They are not presented as a
forecast or as evidence of investable performance. Table~\ref{tab:stats}
lists the statistics produced by the application; percentages are per
observation period.

\begin{table}[!t]
\caption{Bundled Data Set: Sample Statistics}
\label{tab:stats}
\centering
\footnotesize
\begin{tabular}{lrr}
\toprule
Asset & Mean (\%) & Std. dev. (\%)\\
\midrule
Global Equity    & 0.837500 & 1.524243\\
Technology       & 1.266667 & 2.446411\\
Government Bonds & 0.516667 & 0.703202\\
Corporate Bonds  & 0.720833 & 0.358717\\
Gold             & 0.995833 & 1.258529\\
\bottomrule
\end{tabular}
\end{table}

The sample intentionally contains materially different risk-return profiles
and strong positive and negative correlations. This produces a visible
diversification effect and causes several bound constraints to change status
across the frontier. The example is therefore more informative for the active-
set method than a diagonal covariance matrix, while remaining small enough for
manual interpretation.

\subsection{Real-Market Demonstration Data}

In addition to the deterministic synthetic data set, the final application
includes an optimizer-ready real-market sample derived from a public
adjusted-close data set \cite{marketdata} containing SPY, QQQ, IEF, LQD,
and GLD. The data comprise 72 monthly return observations from January 2020
through December 2025. Monthly observations were derived from daily adjusted
closing prices by retaining the final available trading observation of each
calendar month and computing
\begin{equation}
    r_t = \frac{P_t}{P_{t-1}} - 1.
    \label{eq:market-return}
\end{equation}
Both the adjusted-close source subset and the derived monthly-return matrix are
bundled with the project so that the transformation remains reproducible.

For this sample, the estimated periodic mean returns for SPY, QQQ, IEF, LQD,
and GLD are approximately $1.2919\%$, $1.7084\%$, $0.0279\%$, $0.1268\%$, and
$1.5143\%$, respectively. Their corresponding sample standard deviations are
approximately $4.9456\%$, $5.9408\%$, $2.1380\%$, $2.8491\%$, and $4.2563\%$.
The sample is provided as a realistic demonstration of the same optimization
pipeline rather than as evidence of future investment performance.

\subsection{Automated Numerical Tests}

The executable \texttt{PortfolioCoreTests} is registered with CTest. Its
checks are summarized in Table~\ref{tab:tests}. The tests exercise both
successful and rejected inputs; they do not merely verify that the application
runs without an exception.

\begin{table}[!t]
\caption{Automated Verification Scenarios}
\label{tab:tests}
\centering
\scriptsize
\begin{tabular}{p{0.28\columnwidth}p{0.61\columnwidth}}
\toprule
Scenario & Verified property\\
\midrule
CSV and statistics & Dimensions, known sample means, covariance symmetry,
and malformed-number rejection.\\
Two-asset target & Analytical $w_1=w_2=0.5$ solution at the prescribed
return.\\
Diagonal GMV & Analytical weights $(36,9,4)/49$.\\
Active bound & Primal feasibility, a zero-weight asset, and reduced-gradient
KKT conditions.\\
Infeasible target & Rejection above the long-only feasible interval.\\
20-asset diagonal & Dense and sparse feasibility and agreement of all
weights.\\
20-asset correlated & Positive-definite $\mat{B}\mat{B}^{T}+\mat{D}$
covariance, dense/sparse agreement, and an interior optimum.\\
Efficient frontier & 25 feasible, monotonically ordered solutions on the
bundled data set.\\
\bottomrule
\end{tabular}
\end{table}

The KKT test reconstructs the objective gradient with the original covariance
matrix. Multipliers for the budget and return equalities are inferred from two
free assets. Reduced gradients must be near zero for positive weights and must
have the correct sign at active lower bounds. The larger correlated problem
uses $\mat{\Sigma}=\mat{B}\mat{B}^{T}+\mat{D}$ with strictly positive diagonal
$\mat{D}$, which guarantees positive definiteness while retaining nonzero
cross-asset correlations.

A separate backend check solves
\begin{equation}
 \begin{bmatrix}2&1\\1&3\end{bmatrix}
 \begin{bmatrix}x\\y\end{bmatrix}=
 \begin{bmatrix}5\\7\end{bmatrix}
 \label{eq:backend-check}
\end{equation}
with both natID implementations and verifies the known result
$(x,y)=(1.6,1.8)$.

\section{Results and Discussion}

\subsection{Efficient Frontier and Allocations}

With 25 sampled frontier points, the GMV portfolio has expected periodic
return $0.726733\%$, risk $0.134594\%$, and variance
$1.811551\times10^{-6}$. The highest feasible sampled target is
$1.266667\%$, equal to the mean of Technology. Table~\ref{tab:weights}
compares the GMV allocation with the independently optimized $1\%$ target.

\begin{table}[!t]
\caption{Optimal Portfolio Weights}
\label{tab:weights}
\centering
\footnotesize
\begin{tabular}{lrr}
\toprule
Asset & GMV (\%) & Target $1\%$ (\%)\\
\midrule
Global Equity    & 11.483536 & 0.000000\\
Technology       & 0.000000  & 24.316306\\
Government Bonds & 23.552014 & 0.000000\\
Corporate Bonds  & 50.205170 & 22.432726\\
Gold             & 14.759280 & 53.250968\\
\midrule
Total            & 100.000000 & 100.000000\\
\bottomrule
\end{tabular}
\end{table}

The GMV solution assigns its largest weight to Corporate Bonds, the asset with
the smallest individual standard deviation, but it still allocates capital to
three other assets because covariance determines portfolio risk. Technology is
excluded at the GMV point. Raising the target to $1\%$ changes the balance:
Technology supplies the higher expected return, Gold becomes the largest
holding, and the lower-return Global Equity and Government Bonds reach their
active lower bounds. The $1\%$ portfolio attains the target to numerical
precision, has risk $0.281134\%$, variance approximately
$7.90\times10^{-6}$, and converges in six active-set iterations.

Across the 25 sampled points, five active-set transitions create six observed
allocation regions. The sequence begins with Technology fixed at zero near the
GMV portfolio, passes through an interior region, and then successively fixes
Global Equity, Government Bonds, and Corporate Bonds as the target increases.
At the maximum target, Technology has weight one and all other assets are
active. These changes explain the visible piecewise behavior of the weight
paths and changes in curvature of the sampled frontier.

\begin{figure*}[!t]
\centering
\begin{minipage}[t]{0.48\textwidth}
\centering
\imageorplaceholder{Images/efficient_frontier.png}{\linewidth}{2.05in}\\[-1pt]
\small (a)
\end{minipage}
\hfill
\begin{minipage}[t]{0.48\textwidth}
\centering
\imageorplaceholder{Images/portfolio_weights.png}{\linewidth}{2.05in}\\[-1pt]
\small (b)
\end{minipage}
\caption{Portfolio visualizations for the bundled five-asset data set:
(a) the efficient frontier, individual assets, and selected $1\%$ portfolio;
(b) portfolio weights across target returns, with dashed markers at sampled
active-set changes.}
\label{fig:visualizations}
\end{figure*}

Figure~\ref{fig:visualizations}(a) is intended to be read with risk on the
horizontal axis and expected return on the vertical axis. A point above and to
the left is preferable in the mean-variance sense, although the model does not
select a single universally best point because investors may accept different
risk levels. Figure~\ref{fig:visualizations}(b) exposes information that the
frontier alone cannot show: which assets generate each efficient portfolio and
exactly where zero-weight constraints change.

\subsection{Dense and Sparse Backend Agreement}

Table~\ref{tab:backend} reports one comparison captured in the Debug
configuration for the $1\%$ target. Both backends execute the same active-set
logic and differ only in the solution of \eqref{eq:kkt}.

\begin{table}[!t]
\caption{Dense--Sparse Comparison at a $1\%$ Target}
\label{tab:backend}
\centering
\footnotesize
\begin{tabular}{lrr}
\toprule
Quantity & Dense & Sparse\\
\midrule
Iterations & 6 & 6\\
Active constraints & 2 & 2\\
Single-run time (ms) & 0.6823 & 0.3571\\
Budget residual & $0$ & $0$\\
Return residual & $1.735\!\times\!10^{-18}$ &
$1.735\!\times\!10^{-18}$\\
\bottomrule
\multicolumn{3}{l}{\scriptsize Max. weight difference:
$5.551\times10^{-17}$}\\
\multicolumn{3}{l}{\scriptsize Variance difference:
$5.082\times10^{-21}$}\\
\bottomrule
\end{tabular}
\end{table}

The active sets are identical, and the maximum weight difference is close to
double-precision roundoff. This is stronger evidence than matching objective
values alone because two implementations independently reproduce the entire
allocation. The comparison tolerance used by the interface is $10^{-7}$.
The measured times are displayed for diagnostic interest only. They are single
runs from a Debug build, include solver setup, and do not control processor
state or repeat count; they therefore do not establish a general performance
advantage for either storage scheme.

\begin{figure}[!t]
\centering
\imageorplaceholder{Images/backend_comparison.png}{\columnwidth}{1.65in}
\caption{Portfolio Explorer showing the selected portfolio together with
numerical optimization details and the dense--sparse natID backend comparison.}
\label{fig:comparison}
\end{figure}

The sparse representation becomes conceptually more relevant as the number of
assets and active bounds grows, because the lower KKT blocks contain many
structural zeros. However, the covariance block estimated from ordinary asset
returns is generally dense. Consequently, a portfolio KKT matrix is not fully
sparse, and a performance claim requires controlled experiments over multiple
problem sizes. The project deliberately limits its conclusion to numerical
agreement and correct use of both professor-provided matrix implementations.

\section{Limitations}

The numerical results must be interpreted within the model assumptions.
Historical or synthetic sample means are uncertain estimates of future return,
and sample covariance can be poorly conditioned when the observation count is
small relative to the number of assets. A long-only restriction avoids
leveraged and short positions but does not model upper bounds, sector limits,
turnover, transaction costs, taxes, or a risk-free asset. Variance penalizes
positive and negative deviations symmetrically and does not directly represent
tail loss. The model is single-period and static; it does not rebalance through
time.

The bundled data are synthetic and were chosen for deterministic verification
and clear visualization. They demonstrate the numerical method, not a trading
strategy. The bundled real-market sample demonstrates that the same pipeline can operate
on historical market observations, but it remains an in-sample demonstration
rather than an investment study. A practical deployment would additionally
require explicit frequency and currency normalization, systematic treatment of
missing values, out-of-sample evaluation, and sensitivity analysis. From a numerical
perspective, larger or nearly dependent data sets may also require stronger
covariance estimation, scaling, or specialized symmetric-indefinite
factorization. These extensions are intentionally outside the scope of the
implemented course project.

\section{Conclusion}

The project implements the proposed mean-variance model as a complete and
verifiable numerical application. It estimates return statistics from a CSV
matrix, formulates the long-only target-return problem as a convex quadratic
program, and solves it through a feasible primal active-set method. Explicit
working-set updates make zero-weight constraints observable rather than hiding
them behind a library optimizer. Repeated solutions construct the efficient
frontier and show how allocations change as additional return is required.

The implementation also demonstrates a separation between optimization logic
and linear algebra. Dense and sparse natID backends solve the same KKT systems
and agree to near machine precision on the bundled example, while automated
tests extend verification to analytical and 20-asset cases. The resulting
frontier, weight paths, residuals, active constraints, and CSV export provide
multiple independent views of correctness. Within its stated long-only,
single-period assumptions, the application fulfills the intended numerical
optimization objectives and provides a clear basis for explaining both the
mathematics and the implementation.

\balance
\begin{thebibliography}{00}

\bibitem{markowitz1952}
H. Markowitz, ``Portfolio selection,'' \emph{The Journal of Finance},
vol. 7, no. 1, pp. 77--91, Mar. 1952, doi: 10.2307/2975974.

\bibitem{nocedal2006}
J. Nocedal and S. J. Wright, \emph{Numerical Optimization}, 2nd ed.
New York, NY, USA: Springer, 2006.

\bibitem{goldfarb1983}
D. Goldfarb and A. Idnani, ``A numerically stable dual method for solving
strictly convex quadratic programs,'' \emph{Mathematical Programming},
vol. 27, pp. 1--33, 1983, doi: 10.1007/BF02591962.

\bibitem{golub2013}
G. H. Golub and C. F. Van Loan, \emph{Matrix Computations}, 4th ed.
Baltimore, MD, USA: Johns Hopkins University Press, 2013.

\bibitem{davis2006}
T. A. Davis, \emph{Direct Methods for Sparse Linear Systems}.
Philadelphia, PA, USA: Society for Industrial and Applied Mathematics, 2006.

\bibitem{marketdata}
M. Sahni, ``marketdata,'' GitHub repository. [Online]. Available:
\url{https://github.com/manisahni/marketdata}. Accessed: Sep. 18, 2026.

\bibitem{natid}
I. D\v{z}afi\'{c}, ``natID: Native Interface Design for macOS, Linux, and
Windows,'' GitHub repository. [Online]. Available:
\url{https://github.com/idzafic/natID}. Accessed: Sep. 18, 2026.

\end{thebibliography}

\end{document}
